\documentclass[10pt,a4paper]{amsart}
\usepackage[margin=19mm]{geometry}
\usepackage{amsmath,amssymb,mathtools}
\usepackage[scr=rsfs,cal=euler]{mathalfa}
\usepackage{xfrac}
\usepackage[unicode,hidelinks,bookmarks=false]{hyperref}
\newcommand{\ie}{i.e.\ }
\newcommand{\cf}{cf.\ }
\newcommand{\resp}{respectively\ }
\newcommand{\Subsection}{\subsection}
\providecommand{\nobreakdash}{\nobreak\hbox{-}}
\setlength{\emergencystretch}{2em}
\multlinegap0pt
% ---------------------------------------------------------------------------
% Prelude: the paper's own mathematical macros, list and theorem environments,
% transcribed from its own preamble (cached-originals/source0.tex lines 74--185).
% No mathematics is changed; the AMS amsart wrapper is editorial formatting only.
% ---------------------------------------------------------------------------
\newcommand{\TT}{\mathbb{T}}
\newcommand{\NN}{\mathbb{N}}
\newcommand{\ZZ}{\mathbb{Z}}
\newcommand{\KK}{\mathbb{K}}
\newcommand{\QQ}{\mathbb{Q}}
\newcommand{\RR}{\mathbb{R}}
\newcommand{\CC}{\mathbb{C}}
\newcommand{\PP}{\mathbb{P}}
\newcommand{\BBA}{\mathbb{A}}
\newcommand{\FB}{\mathbf{B}}
\newcommand{\FD}{\mathbf{D}}
\newcommand{\fd}{\mathbf{d}}
\newcommand{\frv}{\mathfrak{v}}
\DeclareMathOperator{\rk}{rk}
\DeclareMathOperator{\Dr}{\rm Dr}
\DeclareMathOperator{\proj}{Proj}
\DeclareMathOperator{\Proj}{Proj}
\DeclareMathOperator{\Hom}{Hom}
\DeclareMathOperator{\Spec}{Spec}
\DeclareMathOperator{\Ima}{Im}
\DeclareMathOperator{\Trop}{Trop}
\newcommand{\MC}{\mathcal{C}}
\newcommand{\MI}{\mathcal{I}}
\newcommand{\MH}{\mathcal{H}}
\newcommand{\MX}{\mathcal{X}}
\newcommand{\MY}{\mathcal{Y}}
\newcommand{\ML}{\mathcal{L}}
\newcommand{\MM}{\mathcal{M}}
\newcommand{\MO}{\mathcal{O}}
\newcommand{\MP}{\mathcal{P}}
\newcommand{\MF}{\mathcal{F}}
\newcommand{\MS}{\mathcal{S}}
\newcommand{\mC}{\mathcal{C}}
\newcommand{\mY}{\mathcal{Y}}
\newcommand{\rank}{\textrm{Rank}}
\newcommand{\init}{\textrm{in}_{\prec}}
\newcommand{\Fl}{\mathscr{F}\!\ell}
\newcommand{\coker}{\mathrm{coker}}
\newcommand{\bigzero}{\mbox{\normalfont\Large\bfseries 0}}
\newcommand{\Bsig}{\mathcal{B}^\sigma}
\newcommand{\Btau}{\mathcal{B}^\tau}
\newcommand{\Bwo}{\mathcal{B}^{w_0}}
\DeclareMathOperator{\argmin}{argmin}
\DeclareMathOperator{\Gr}{Gr}
\DeclareMathOperator{\conv}{Conv}
\DeclareMathOperator{\Ker}{Ker}
\DeclareMathOperator{\vol}{Vol}
\DeclareMathOperator{\sgn}{sgn}
\DeclareMathOperator{\Sym}{Sym}
\newcommand\Myperm[2][^n]{\prescript{#1\mkern-2.5mu}{}P_{#2}}

\theoremstyle{plain}
\newtheorem{theorem}{Theorem}
\newtheorem{lemma}[theorem]{Lemma}
\newtheorem{corollary}[theorem]{Corollary}
\newtheorem{prop}[theorem]{Proposition}
\newtheorem{conjecture}[theorem]{Conjecture}
\newtheorem{question}{Question}
\newtheorem*{claim}{Claim}
\newtheorem{procedure}{Procedure}

\theoremstyle{definition}
\newtheorem{definition}[theorem]{Definition}
\newtheorem{remark}[theorem]{Remark}
\newtheorem{case}{Case}
\newtheorem{subcase}{Case}
\numberwithin{subcase}{case}
\newtheorem*{notation}{Notation}
\newtheorem{example}[theorem]{Example}
\begin{document}
\title{Generalized metric tree arrangements and Dressians}
\author{Ayush Kumar Tewari}
\address{Johann Radon Institute for Computational and Applied Mathematics (RICAM)\\
Austrian Academy of Sciences, Linz, Austria}
\email{tropical.tewari@gmail.com, ayushkumar.tewari@ricam.oeaw.ac.at}
\begin{abstract}
Metric trees and metric tree arrangements index cones in the polyhedral fan structure in the Dressian $\Dr(2,n)$ and $\Dr(3,n)$ respectively. We introduce the notion of generalized metric tree arrangements which parameterize points in $\Dr(k,n)$ and extend previously known results to $\Dr(k,n)$ along with providing explicit examples of these generalized metric tree arrangements. We study the adjacency of cones in the positive Dressian $\Dr_{>0}(3, n)$ and introduce generalized Whitehead moves which provide a condition for adjacency of maximal cones in $\Dr_{>0}(3,n)$ in terms of the associated metric tree arrangements. 
\end{abstract}

\maketitle
\noindent\emph{Source, version and licence.} \emph{Generalized metric tree arrangements and Dressians}, by Ayush Kumar Tewari. The e-print used for this editable unit is \textbf{version v2} of arXiv:2210.06868, https://arxiv.org/abs/2210.06868v2, whose material is distributed under the Creative Commons Attribution 4.0 International licence, http://creativecommons.org/licenses/by/4.0/, as recorded in the arXiv licence metadata and shown on that abstract page. Version v3 of arXiv:2210.06868 is an arXiv administrative withdrawal: the v3 e-print and the v3 PDF both return HTTP 404 and the abstract submission history lists v3 as withdrawn with a 1 KB entry, so no v3 material exists and none is used here; v2 is the newest intact version of the paper and the whole of the body below is taken from it. Full rights notice: licenses/arxiv-2210.06868-CC-BY-4.0.txt.
\par\smallskip
\noindent\textit{Editorial scope.} Counted unit: original Theorem 23 of the article, the statement carried by the source environment \texttt{theorem} with source label \texttt{prop:gen\_tree\_dressian}, cached-originals/source0.tex lines 656--658: ``As sets, the collection of weights $\{w_T \colon T \in \TT(k,n) \}$ is equal to the Dressian $\Dr(k,n)$.'' It is delivered together with the author's own complete proof as the single primary proof body, source0.tex lines 660--671, that is the \texttt{proof} environment which immediately follows the statement, which the author opens with ``Let $T \in \TT(k,n)$ be a tree arrangement'' and which establishes both inclusions of the set equality before it closes. No proof marker is added and no mathematics is written, repaired or re-derived; the printed slips of the source are kept literal. Necessary complete context is retained uncounted: source0.tex lines 202--204, the author's own abstract; lines 228--246, the section Dressian from its heading through the definition of the Grassmannian, of the tropical Grassmannian and of the Dressian $\Dr(k,n)$ as the tropical prevariety of the three-term Plucker relations; lines 267--269, the definition of a phylogenetic tree, of its tree distance vector and of the space of phylogenetic trees; lines 401--418, the definition of a metric tree arrangement and of an abstract tree arrangement; lines 568--581, the section Generalized metric tree arrangements for $\Dr(k,n)$ from its heading through its introductory sentence and Lemma \texttt{lem:mat\_sub\_k,n} (source lines 573--575) with its complete author proof (source lines 577--581); lines 636--654, the introductory sentence, the definition of a generalized metric tree arrangement with its compatibility equation and the paragraph which defines the induced weight vector $w_T$; lines 673--692, the alternative description through the map $\pi$, the induced fan structure on the space of weights and the notion of a cherry; lines 694--701, the corollary that each regular matroidal subdivision of $\Delta(k,n)$ supports a generalized metric tree arrangement, with its complete author proof; lines 707--717, the equivalence discussion and the concluding corollary that the equivalence classes of generalized metric tree arrangements are in bijection with the regular matroidal subdivisions of $\Delta(k,n)$, with its complete author proof; lines 825--829, the corollary condensing the section; and line 831, the author's acknowledgement. The original shared theorem-like counter is preserved through explicit counter settings, so the counted statement is printed as Theorem 23 exactly as in the article and the retained Lemma, Definitions, Remark and Corollaries keep their original numbers. All references inside the retained spans resolve to labels defined inside the retained spans. Omitted from this editable unit and therefore absent from it: source0.tex lines 1--201 and 205--227, that is the article's own class and package preamble, its title and affiliation block and its table of contents; lines 247--266, 270--400, 419--567 and 582--635, that is the remainder of the section Dressian, the whole of the section Adjacency of maximal cones in positive Dressian including its figure-dependent paragraph, and the examples for $\Delta(4,8)$ and $\Delta(4,6)$; lines 806--824 and 826, that is the closing remarks and the open questions; lines 832--1166, that is the bibliography and the whole appendix with its tables of cherries and adjacent cones. All 22 figure environments of the e-print, at source0.tex lines 289--294, 393--398, 547--552, 622--627, 772--777, 799--804, 868--873, 912--917, 933--938, 954--959, 975--980, 996--1001, 1017--1022, 1057--1062, 1075--1080, 1084--1089, 1093--1098, 1102--1107, 1111--1116, 1120--1125, 1129--1134 and 1139--1144, are omitted, together with every one of the 23 artwork PDF files shipped in the e-print. No \texttt{\textbackslash includegraphics} command and no figure environment occurs anywhere in this delivered file, so no artwork is redistributed. The bibliography of this unit is restricted to the five entries actually cited in the retained spans, transcribed from the author's own \texttt{main.bbl}; they keep the author's own numeric labels 7, 8, 11, 13 and 15, and the author's own \texttt{\textbackslash bibliographystyle} and \texttt{\textbackslash bibliography} lines and the remaining eleven entries of that file are omitted. Every declared body of this unit is an exact, unmodified span of source0.tex: no delivered body is a bounded passage comparison, \texttt{omit} was not used anywhere, and consequently no fidelity receipt is asserted in delivery-evidence.json. The AMS amsart wrapper, the named format packages of the pinned TeX Live runtime and the editorial source, licence and scope paragraphs are editorial formatting only.
\par\smallskip
\section{Dressian}\label{sec:Section_2}

%\subsection{The structure of the Dressian Dr(2,n) }

We recall some basic definitions of Grassmannian, tropical Grassmannian and Dressian \cite{maclagan2021}.
The Grassmannian $\Gr(k,n)$ is a smooth projective variety defined by the Pl\"ucker ideal $I_{k,n}$,

\[ I _{k,n} = \langle  \mathcal{P}_{I,J} : I,J \subseteq [n], |I| = k-1, |J| = k+1 \rangle  \]

whose generators are the quadratic Pl\"ucker relations,

\[ \mathcal{P}_{I,J} = \sum_{j \in J} \text{sgn}(j;I,J) \cdot p_{I \cup j} \cdot p_{J \setminus j }  \]

The Grassmannian parameterizes $k$ dimensional subspaces of the vector space $\mathbb{K}^{n}$, where $\mathbb{K}$ is the underlying field. The \emph{tropical Grassmannian} $\Trop(\Gr(k,n))$ is the tropical variety defined by the Pl\"ucker ideal $I_{k,n}$. Among the generators of the ideal $I_{k,n}$ are the \emph{three term Pl\"ucker} relations which are of the form

\[ p_{Apq}p_{Ars} - p_{Apr}p_{Aqs} + p_{Aps}p_{Aqr}\]


where $A \in \binom{[n]}{k-2}$ and $p,q,r,s \in [n] \setminus A$ are pairwise distinct. The three-term Pl\"ucker relations generate the image of the ideal $I_{k,n}$ in the Laurent polynomial ring $\mathbb{K}[p_{A}^{\pm 1}]$, where  $A \in $ $\binom{[n]}{k-2}$. The \emph{Dressian} $\Dr(k,n)$ is the tropical prevariety defined by all the three-term Pl\"ucker relations. The Grassmannian and Dressian have a pure rational polyhedral fan structure in $\mathbb{R}^{{\binom{n}{k}} - 1 } \cong  \mathbb{R}^{\binom{n}{k}} / \mathbb{R}1$.

We firstly discuss the case of $\Dr(2,n)$. $\Trop(\Gr(2,n))$ enjoys a strong connection with the space of phylogenetic trees which we recall here.

A \emph{phylogenetic} tree $T$ is a tree with $n$ labeled leaves and no vertices of degree two. A phylogenetic tree $T$ is \textit{trivalent} if for each vertex $v \in V(T)$ either $\deg(v) = 3$ or $\deg(v) = 1$. To each such phylogenetic tree we can attach metric data in the following way. We assign a length $l_{e} \in \mathbb{R}$ to each edge of $T$. Between any two leaves $i$ and $j$ of $T$, there is a unique path in $T$, and the distance between the leaves can be obtained as sum of $l_{e}$ along this path. This gives us the \emph{tree distance} vector $\delta = (\delta_{ij}) \in \mathbb{R}^{\binom{n}{2}}$. When all $l_{e}$ are nonnegative, the tree distance gives a finite metric. Finite metric spaces that arise from such metric tree are called \emph{tree metric}. The set of all tree distances is known as the space of phylogenetic trees on $n$ leaves and is denoted as $\mathcal{T}_{n}$ \cite{maclagan2021}.

\setcounter{theorem}{8}
\begin{definition}

Let $n \geq 4$ and we consider a $n$-tuple of metric trees $T = (T_{1}, \hdots , T_{n})$, each with an associated metric $\delta_{i}$. We call the $n$- tuple $T$ of metric trees a \emph{metric tree arrangement} if, 

\[ \delta_{i}(j,k) = \delta_{j}(i,k) =\delta_{k}(i,j) \]

for all $i,j,k \in [n]$ pairwise distinct \cite{herrmann2009draw}.

Moreover, considering trees $T_{i}$ without metrics, but
with leaves still labeled by $[n] \setminus {i}$, we say that $T$ is an \emph{abstract tree arrangement} if,

\begin{itemize}
    \item $n =4$,
    \item $n = 5$, and $T$ is the set of quartets of a tree with five leaves,
    \item $n \geq 6$, and $(T_{1} \setminus i, \hdots , T_{i-1} \setminus i, T_{i+1} \setminus i, \hdots , T_{n} \setminus i)$ is an arrangement of $n - 1$ trees for each $i \in [n]$ \cite{herrmann2009draw}.
\end{itemize}

\end{definition}

\setcounter{theorem}{18}
\section{Generalized metric tree arrangements for Dr(k,n)}\label{sec:Section_4}

We begin this section with a discussion on the structure of Dressian  
$\Dr(k,n)$ ($k \geq 4)$. Firstly, we state Lemma \ref{lem:mat_sub_k,n} which is a generalization of \cite[Lemma 4.3]{herrmann2009draw},

\begin{lemma}\label{lem:mat_sub_k,n}
Each matroid subdivision $\Sigma$ of $\Delta(k,n)$ defines an abstract tree arrangement $T(\Sigma)$ of $\Myperm[n]{k-2}$ trees, with each tree having $n-(k-2)$ leaves. Moreover, if $\Sigma$ is regular then $T(\Sigma)$ supports metric tree arrangement.
\end{lemma}

\begin{proof}
Each of the $n$ contraction facets of $\Delta(k,n)$ are  isomorphic to $\Delta(k-1,n-1)$. Recursively using this relation we get that $\Sigma$ induces matroidal subdivisions on $(k-2)! \cdot {\binom{n}{k-2}} = \Myperm[n]{k-2}$ copies of $\Delta(2,n-(k-2))$. But matroidal subdivisions of $\Delta(2,n-(k-2))$ are generated by the compatible system of splits, and such subdivisions are dual to trees. Hence, $\Sigma$ gives rise to a tree arrangement.

Let $\Sigma$ be regular with $\lambda$ as a weight function. We now invoke arguments, similar to what has been used in proof of Lemma 4.2 in \cite{herrmann2009draw}. By suitable rescaling, we can assume that $\lambda$ has values between $1$ and $2$. We restrict $\lambda$ to $\Delta(2,n-(k-2))$ to obtain the induced subdivisions. But a weight function on $\Delta(2,n-(k-2))$ which takes values between $1$ and $2$ is a metric. Since, the induced subdivisions on $\Delta(2,n-(k-2))$ are also regular matroidal subdivisions, they are dual to metric trees on $n-(k-2)$ leaves. The Split Decomposition theorem \cite[Theorem 2.11]{herrmann2008splitting} provides the lengths on all edges of these trees.
\end{proof}

\setcounter{theorem}{20}
We now describe an extension of the definition of metric tree arrangement \cite{herrmann2009draw}. 

\begin{definition}\label{def:gen_metric_tr_arr}
A \textit{generalized metric tree arrangement for $\Dr(k,n)$} is a set $T = \{T_I \colon I \subseteq [n], \ |I| = k-2 \}$ of metric trees such that the leaves of $T_I$ are labeled by the elements of $[n] \backslash I$ and $T$ satisfies the following compatibility condition: for any pair of metric trees $T_I$ and $T_{I'}$ in $T$; any pair of leaves $i, j$ of $T_I$; and any pair of leaves $i', j'$ of $T_{I'}$, if $I \cup \{i,j\} = I' \cup \{i', j'\}$ then we have

\begin{equation}\label{eq:comptability}
 \delta_I(i,j) = \delta_{I'}(i', j')
\end{equation}

where $\delta_I$ is the metric on $T_I$ and $\delta_{I'}$ is the metric on $T_{I'}$. We define $\TT(k,n)$ to be the collection of all generalized metric tree arrangements for $\Dr(k,n)$.
\end{definition}


\begin{remark}
It is evident from the definition that a metric tree arrangement is also a generalized metric tree arrangement and corresponds to the case when $k=3$.
\end{remark}


Each generalised metric tree $T$ induces a weight vector $w_T$ on the Pl\"ucker variables given by $w_T(P_K) = \delta_J(i,j)$ where: $K = J \cup \{i,j\}$; $\delta_J$ is the metric on the tree $T_J$; and $i,j$ are leaves of the tree $T_J$. Note that the definition of the metric tree arrangement above ensures that $w_T$ is a well-defined weight on the Pl\"ucker variables. We establish the relation between the weight vector $w_{T}$ and the Pl\"ucker fan structure of $\Dr(k,n)$.

\setcounter{theorem}{22}
\begin{theorem}\label{prop:gen_tree_dressian}
As sets, the collection of weights $\{w_T \colon T \in \TT(k,n) \}$ is equal to the Dressian $\Dr(k,n)$.
\end{theorem}

\begin{proof}
Let $T \in \TT(k,n)$ be a tree arrangement. Fix a $3$-term tropical Pl\"ucker relation for $\Dr(k,n)$. Note that this relation is of the form
\[
P = 
P_{J \cup ij} \odot P_{J \cup k\ell} \oplus
P_{J \cup ik} \odot P_{J \cup j\ell} \oplus
P_{J \cup i\ell} \odot P_{J \cup jk} 
\]
where $J \subseteq [n]$ has size $k - 2$ and $i,j,k,\ell \in [n] \backslash J$ are distinct. Let $S = \{I \subseteq [n] \colon |I| = k, \ J \subseteq I \}$ be the $k$-subsets containing $J$. By definition, the restriction of $w_T$ to the Pl\"ucker variables $P_I$ where $I \in S$ is induced by a single metric tree $T_J \in T$. By \cite{speyer2004tropical}, the space of weights induced by metric trees with $n-k+2$ leaves equals the Dressian $\Dr(2, n-k+2)$. Therefore the maximum in $P$ is attained by at least two terms. So we have shown that $\{w_T \colon T \in \TT(k,n) \} \subseteq \Dr(k,n)$.

For the other inclusion take any point $w \in \Dr(k,n)$. For each $(k-2)$-subset $J \subseteq [n]$, we construct a metric tree $T_J$ with leaves $[n]\backslash J$. Consider the restriction of $w$ to the Pl\"ucker variables $P_I$ where $J \subseteq I$. Note that these Pl\"ucker variables are naturally in bijection with the $2$-subsets of $[n]\backslash J$. So, $w$ gives rise to point in $\Dr(2, n-k+2)$ given by $\hat w(P_{ij}) = w(P_{J \cup ij})$. As above, the Dressian $\Dr(2, n-k+2)$ is equal to the space of weights induced by metric trees and so we define $T_J$ to be the metric tree associated to $\hat w$ with leaves labeled by $[n] \backslash J$. It remains to show that the trees satisfy the compatibility condition. However, this is immediate since each metric tree $T_J$ is constructed from the same point $w \in \Dr(k,n)$.
\end{proof}

Alternatively, the result in Theorem \ref{prop:gen_tree_dressian} can be understood as a generalization of the map $\pi$ defined in the proof of Theorem 4.4 in \cite{herrmann2009draw}. We describe this perspective here. Let $\pi$ be the map from $[n]^{k} \rightarrow \mathbb{R} \cup \{\infty\} $

\[ \pi(i_{1}, i_{2}, \hdots , i_{k}) = \begin{cases}
  \delta_{I_{1}}(i_{1},i_{2}) =  \delta_{I_{1}}(i_{3},i_{4}) \hdots = \delta_{I_{m}}(i_{k-1},i_{k}) &  i_{1}, i_{2}, \hdots , i_{k} \> \text{are pairwise distinct} \\    
  0 & \text{otherwise}    
\end{cases} \]

where $I_{1}, \hdots I_{m}$ are compatible sets of size $k-2$ satisfying compatibility conditions described in Definition \ref{def:gen_metric_tr_arr}. With the proof of Theorem \ref{prop:gen_tree_dressian} we see that for a compatible set $I \in$ $\binom{[n]}{k}$, the minimum of 

\[ \text{min} \{ \pi_{Ipq} + \pi_{Irs} , \pi_{Ipr} + \pi_{Isq} , \pi_{Ips} + \pi_{Iqr} \} \]

is attained twice, where $p,q,r,s \in [n] \setminus I$, as this minimum corresponds to the weight induced by a metric tree on $n-k+2$ leaves. Hence, we can conclude that the restriction of the map $\pi$ to increasing tuple $i_{1} < i_{2} < \hdots < i_{k}$ is a finite tropical Pl\"ucker vector and is an element of $\Dr(k,n)$.


There is a natural fan structure on the space of weights induced by generalized metric tree arrangements. Two weights $w_T$ and $w_{T'}$ lie in the interior of the same cone if and only if $T$ and $T'$ have the same underlying set of labeled trees. The proof above naturally extends and shows that the Pl\"ucker fan structure of the Dressian coincides with the fan structure on the space of weights induced by generalized metric tree arrangements. Also, as a consequence of Pl\"ucker fan structure and secondary fan structure coinciding for the Dressian \cite{olarte2019local}, the fan structure on the space of weights induced by generalized metric tree arrangements coincides with the secondary fan structure of Dressian. 

The notion of a \textit{cherry} of an arrangement naturally extends to the generalized setting. A \textit{cherry} of an arrangement $T \in \TT(k,n)$ is a $k$-subset $I \in [n]$ such that if $I = J \cup \{i,j\}$ is any partition then $i,j$ is a cherry in tree $T_J$.
\bigskip

With Theorem \ref{prop:gen_tree_dressian} we see that a generalized metric tree arrangement corresponds to a regular matroidal subdivision of $\Delta(k,n)$. We see this when we consider the collection of weights $\{w_T \colon T \in \TT(k,n) \}$ as a height function on the Pl\"ucker coordinates. Using the fact that the Pl\"ucker fan structure and secondary fan structure coincide for the Dressian, we obtain the following result,

\begin{corollary}\label{lem:sub_to_gen_trees}
Each regular matroidal subdivision $\Sigma$ of $\Delta(k,n)$ supports a generalized metric tree arrangement.  
\end{corollary}


\begin{proof}
Let us consider the weight function for the regular matroidal subdivision $\Sigma$ to be $w_{\Sigma}$ and $w_{\Sigma} \in \Dr(k,n)$. Since the Pl\"ucker fan structure and secondary fan structure coincide for the Dressian, therefore by Theorem \ref{prop:gen_tree_dressian}, we can conclude that there exists a generalized metric tree arrangement $T(k,n)$ such that it induces $w_{\Sigma}$.
\end{proof}

We know that in \cite{herrmann2009draw} Lemma 4.2 and Theorem 4.5 help establish a bijection between equivalence classes of \emph{equivalent} metric tree arrangements of $n$ trees and regular matroidal subdivisions of $\Delta(3,n)$. We now try to establish this equivalence for generalized metric tree arrangements. We already see a mapping from one side, i.e., by Theorem \ref{prop:gen_tree_dressian} we see that each generalized metric tree arrangement corresponds to a regular matroidal subdivision of $\Delta(k,n)$. We now have to show a mapping in the reverse direction. 

We refer to a generalized metric tree arrangement as an \emph{abstract} generalized metric tree arrangement if we consider the tree arrangement without the metric. Two generalized metric tree arrangements are \emph{equivalent} if they induce the same abstract tree arrangement. With these results, we can now state the following result which in essence is a generalization of Theorem 4.4 in \cite{herrmann2009draw}.

\begin{corollary}
The equivalence classes of arrangements of $\binom{n}{k-2}$ generalized metric tree are in bijection with the regular matroid subdivisions of $\Delta(k,n)$.
\end{corollary}

\begin{proof}
With Theorem \ref{prop:gen_tree_dressian}, we see that a generalized metric tree arrangement defines a point in the Dressian $\Dr(k,n)$. We observe that equivalent generalized metric tree arrangements define points in the $\Dr(k,n)$ which lie in the same cone in the Pl\"ucker fan. Considering this weight vector as a lifting function provides us the corresponding regular matroid subdivision. This provides us with a map from equivalence classes of generalized metric tree arrangements to regular matroidal subdivisions. For the map from the other side we consider a regular matroidal subdivision of $\Dr(k,n)$ and by Corollary \ref{lem:sub_to_gen_trees} we see that this subdivision corresponds to a generalized metric tree arrangement. By \cite{olarte2019local}, we know that the secondary fan structure and the Pl\"ucker fan structure coincides for $\Dr(k,n)$. Hence, this gives us the required bijection.
\end{proof}

\setcounter{theorem}{28}
We condense some of the findings in this section with this statement,

\begin{corollary}\label{cor:final}
Each maximal cone of the Dressian $\Dr(k,n), n \geq 4, k \geq 2$ supports an arrangement of $\binom{n}{k-2}$ generalized metric trees.
\end{corollary}

\paragraph{\bf Acknowledgment} I am grateful to my advisor Prof Isabel Van Driessche and also to Prof Tom De Medts for their support during the preparation of this article. I would also like to thank Prof Michael Joswig for going through earlier drafts of this work. I thank Dr. Ahmed Umer Ashraf for the discussions which helped in coming up with the examples in the text. A.~K.~Tewari has been supported by the UGent BOF grant BOF/STA/201909/038 and Ugent BOF grant BOFFF12015001101.

\begin{thebibliography}{99}
\bibitem[7]{herrmann2009draw} S.~Herrmann, A.~Jensen, M.~Joswig, and B.~Sturmfels. How to draw tropical planes. {\em The Electronic Journal of Combinatorics}, 16(2):R6, 2009.
\bibitem[8]{herrmann2008splitting} S.~Herrmann and M.~Joswig. Splitting polytopes. {\em M{\"u}nster Journal of Mathematics}, 2008.
\bibitem[11]{maclagan2021} D.~Maclagan and B.~Sturmfels. {\em Introduction to tropical geometry}, volume 161. American Mathematical Society, 2021.
\bibitem[13]{olarte2019local} J.~A. Olarte, M.~Panizzut, and B.~Schr{\"o}ter. On local {D}ressians of matroids. In {\em Algebraic and Geometric Combinatorics on Lattice Polytopes: Proceedings of the Summer Workshop on Lattice Polytopes}, pages 309--329. World Scientific, 2019.
\bibitem[15]{speyer2004tropical} D.~Speyer and B.~Sturmfels. The tropical {G}rassmannian. {\em Advances in Geometry}, 4(3):389--411, 2004.
\end{thebibliography}
\end{document}
